
This work investigates whether the calibration quality of large language model (LLM) code verifiers, measured via Expected Calibration Error (ECE), varies systematically with problem difficulty. Using a self-contained methodology on EvalPlus (542 HumanEval+ and MBPP+ problems), problems are stratified into hard and easy tiers by each model's own k=5 pass@1 distribution, and the metric DELTA\_ECE = ECE(hard) - ECE(easy) is computed with bootstrap confidence intervals. Three base models at 7-8B scale with distinct training regimes are evaluated: DeepSeek-Coder-6.7B (code-specialized pre-training), CodeLlama-7B (code-adapted), and Llama3-8B (general-purpose). The original hypothesis that DELTA\_ECE > 0 universally (i.e., hard problems are more miscalibrated) is refuted: only 1/3 models satisfy this criterion. Instead, the results reveal architecture-dependent calibration behavior. DeepSeek-Coder shows DELTA\_ECE = +0.298 (95\% CI: [0.285, 0.312]), confirming the expected pattern. CodeLlama shows DELTA\_ECE = -0.249 (95\% CI: [-0.259, -0.239]), an inversion where easy problems exhibit greater miscalibration. Llama3 shows DELTA\_ECE approximately 0 (95\% CI includes zero, p = 0.256). Global temperature scaling neither corrects nor unifies these patterns. These findings indicate that P(True) calibration behavior under difficulty stratification is an architectural property rather than a universal phenomenon.
